% Source-aligned English for printed p.195.
\clearpage
\SourcePage{195}

\LayerTarget{tohoku:proof:remark-4.4.1:1}
Since $\mathbf{Ext}_{\mathbf Z(G)}^n(\mathbf Z,A)$ is the sheaf associated to
the presheaf of $\mathrm{Ext}_{\mathbf Z(G)}^n(U;\mathbf Z,A)$, which vanishes
by the foregoing (since $A|U$ is also injective), our assertion is proved.

\LayerTarget{tohoku:remark:4.4.1:2}
\noindent 2. Suppose that $A$ is the \textit{constant} sheaf defined by a
$G$-module $M$. We shall then of course write $H^p(X;G,M)$ and
$\mathbf H^p(G,M)$ in place of $H^p(X;G,A)$ and $\mathbf H^p(G,A)$. If $G$
acts trivially on $M$, then $H^p(X;G,M)$ is therefore computed completely by
formula 4.4.3, or by formula 4.4.2 if $M$ is a vector space over a field $k$.
In the latter case, if moreover $X$ is an HLC space, we even obtain:

\LayerTarget{tohoku:formula:bilinear-equivariant-cohomology}
\[
H^*(X;G,M)=\text{the space of bilinear maps from }
H_*(G,k)\times H_*(X,k)\text{ to }M.
\]

One should note that, in general, $\mathbf H^p(G,M)$ is not the constant sheaf
defined by $H^p(G,M)$; cf. Remark 2 of 4.3.

\LayerTarget{tohoku:chapter:V}
\bigskip
\begin{center}
{\large\bfseries Chapter V\quad Cohomological Study of Spaces with Operators}
\end{center}

\LayerTarget{tohoku:section:5.1}
\noindent\textbf{5.1. Generalities on $G$-sheaves.}

\LayerTarget{tohoku:setup:g-space-and-orbit-space}
Throughout this chapter we consider a space $X$ on which a group $G$ acts
(on the left, to fix ideas). The operation defined by $g\in G$ will be denoted
$x\mapsto g\mathbin{\cdot}x$. We do not require the action of $G$ to be
faithful; in particular, we shall also have to consider the case in which $G$
acts trivially. We write $X(G)$ for $X$ endowed with the additional structure
defined by the operations of $G$. The orbit space $X/G$ will be denoted by
$Y$ and endowed with the quotient topology. The canonical map from $X$ onto
$Y$ will be denoted by $f$; it is a continuous open map. We regard $Y$ as
endowed with the operator group $G$ acting trivially, and hence write $Y(G)$
to recall this additional structure on $Y$.

\LayerTarget{tohoku:definition:g-sheaf}
By a \textit{$G$-sheaf} on $X=X(G)$ we mean a sheaf (of sets) $A$ on $X$ on
which $G$ acts compatibly with its operations on $X$. To give a meaning to
this definition, one may for example regard $A$ as an étalé space over $X$
(cf. 3.1); we shall not dwell on this. One defines in the same way the notions
of a \textit{$G$-sheaf of groups, of rings}, and of an \textit{abelian
$G$-sheaf} (= a $G$-sheaf of abelian groups), etc.

\LayerTarget{tohoku:example:g-sheaves}
Figuratively, one may say that if $X$ is endowed with certain structures and a
sheaf $A$ on $X$ is defined in structural terms, then $A$ is naturally a
$G$-sheaf whenever the operations of $G$ on $X$ are automorphisms. Thus a
constant sheaf can always be regarded as a $G$-sheaf (a ``trivial
$G$-sheaf''); likewise for the sheaf of germs of arbitrary, respectively
continuous, maps from $X$ to a given set, and for the sheaf of germs of
holomorphic functions when $X$ is a holomorphic variety and the operations of
$G$ are automorphisms of the variety $X$, etc.

\LayerTarget{tohoku:definition:g-homomorphism}
A \textit{$G$-homomorphism} from one sheaf to another is a sheaf homomorphism
that commutes with the operations of $G$. If, for example, the sheaves under
consideration are sheaves of groups, it is understood as usual that the
homomorphism respects that structure.

\LayerTarget{tohoku:definition:category-g-sheaves}
With this notion of homomorphism, the $G$-sheaves of sets (respectively the
$G$-sheaves of groups, etc.) form a \textit{category} (cf. 1.1) having the
same properties as the corresponding category without an operator group $G$
(which is, moreover, a special case, corresponding to the case in which $G$
is reduced to its identity element).

\TohokuContinuation{196}
